% !TeX program = xelatex
\documentclass[aspectratio=169, 10pt]{beamer}

% ==============================================================================
% GÓI VÀ CẤU HÌNH HỆ THỐNG
% ==============================================================================
\usepackage{fontspec}
\setsansfont{Arial}
\setmonofont{Consolas}

\usepackage{booktabs}
\usepackage{colortbl}
\setlength{\extrarowheight}{2pt}
\usepackage{array}
\usepackage{tikz}
\usetikzlibrary{shapes.geometric, arrows.meta, positioning, calc}
\usepackage{listings}

% ==============================================================================
% THIẾT LẬP PHONG CÁCH CHUẨN TEMPLATE: NỀN TRẮNG - CHỮ ĐEN - BẢNG MÀU XANH
% ==============================================================================
\useinnertheme{default}
\useoutertheme{default}

% Màu bảng chuẩn PowerPoint theo yêu cầu (Nền xanh, viền trắng)
\definecolor{TableHead}{RGB}{68, 114, 196}      % Xanh dương đậm (#4472C4)
\definecolor{TableRowOdd}{RGB}{207, 213, 234}   % Xanh tím nhạt dòng lẻ (#CFD5EA)
\definecolor{TableRowEven}{RGB}{233, 235, 245}  % Xanh tím rất nhạt dòng chẵn (#E9EBF5)

% Cấu hình màu Beamer thuần phong cách Nền trắng - Chữ đen
\setbeamercolor{background canvas}{bg=white}
\setbeamercolor{normal text}{fg=black}
\setbeamercolor{frametitle}{fg=black, bg=white}
\setbeamercolor{title}{fg=black}
\setbeamercolor{subtitle}{fg=black}
\setbeamercolor{author}{fg=black}
\setbeamercolor{institute}{fg=black}
\setbeamercolor{date}{fg=black}
\setbeamercolor{item}{fg=black}
\setbeamercolor{subitem}{fg=black}
\setbeamercolor{block title}{fg=black, bg=gray!20}
\setbeamercolor{block body}{fg=black, bg=gray!8}
\setbeamercolor{block title alerted}{fg=black, bg=gray!20}
\setbeamercolor{block body alerted}{fg=black, bg=gray!8}

% Căn lề tiêu đề slide thuần đen thanh lịch theo template chuẩn
\setbeamertemplate{frametitle}{
  \vspace{0.25cm}
  \textbf{\large\color{black}\insertframetitle}
  \par\vspace{0.1cm}
  {\color{black}\hrule height 0.6pt}
  \vspace{0.2cm}
}

% Thanh chân trang cực kỳ tối giản, không màu mè
\setbeamertemplate{navigation symbols}{}
\setbeamertemplate{footline}{
  \leavevmode%
  \hbox{%
  \begin{beamercolorbox}[wd=.5\paperwidth,ht=2.5ex,dp=1ex,left]{footline}%
    \hspace*{2ex}\scriptsize\color{gray} Lập trình Đa nền tảng
  \end{beamercolorbox}%
  \begin{beamercolorbox}[wd=.5\paperwidth,ht=2.5ex,dp=1ex,right]{footline}%
    \scriptsize\color{gray} Trang \insertframenumber{} / \inserttotalframenumber{}\hspace*{2ex}
  \end{beamercolorbox}}%
  \vskip3pt%
}

% Cấu hình hiển thị code đơn giản (nền trắng, viền mảnh, chữ đen)
\lstset{
  language=Java,
  basicstyle=\ttfamily\scriptsize\color{black},
  keywordstyle=\bfseries\color{black},
  stringstyle=\color{black},
  commentstyle=\color{gray}\itshape,
  numbers=none,
  showstringspaces=false,
  breaklines=true,
  frame=single,
  rulecolor=\color{gray!40},
  backgroundcolor=\color{white},
  tabsize=2,
  escapeinside={(*@}{@*)},
  xleftmargin=3pt,
  xrightmargin=3pt,
  aboveskip=3pt,
  belowskip=3pt
}

% Thông tin đề tài
\title{\textbf{CHỦ ĐỀ 01: OOP NÂNG CAO \& XỬ LÝ NGOẠI LỆ TẬP TRUNG TRONG DART}\\ \large XÂY DỰNG MODULE QUẢN LÝ DỮ LIỆU AN TOÀN, LOẠI BỎ CRASH VÀ CÁCH LY LỖI}
\subtitle{Môn học: Lập trình Đa nền tảng}
\author{\textbf{SVTH:} 1. Lê Anh Khoa \quad 2. Ngô Quang Thế}
\institute{
  \textbf{TRƯỜNG ĐẠI HỌC BÁCH KHOA}\\
  \textbf{KHOA ĐIỆN TỬ VÀ TRÍ TUỆ NHÂN TẠO}
}
\date{Đà Nẵng, 2026}

\begin{document}

% ==============================================================================
% SLIDE 1: TRANG BÌA (CÂN ĐỐI, NỀN TRẮNG CHỮ ĐEN THEO CHUẨN ĐẠI HỌC)
% ==============================================================================
\begin{frame}[plain]
  \vspace{0.4cm}
  \begin{center}
    {\large \textbf{TRƯỜNG ĐẠI HỌC BÁCH KHOA -- ĐẠI HỌC ĐÀ NẴNG}}\\[0.12cm]
    {\large \textbf{KHOA ĐIỆN TỬ VÀ TRÍ TUỆ NHÂN TẠO}}
  \end{center}

  \vspace{0.45cm}
  \begin{center}
    {\normalsize Môn học: \textbf{Lập trình Đa nền tảng}}
  \end{center}

  \vspace{0.45cm}
  \begin{center}
    {\LARGE \textbf{<<Chủ đề 01: OOP nâng cao \& Xử lý ngoại lệ>>}}\\[0.22cm]
    {\normalsize Module quản lý dữ liệu an toàn, loại bỏ crash và cách ly lỗi trong Dart}
  \end{center}

  \vspace{0.45cm}
  \begin{center}
    \begin{tabular}{rl}
      \textbf{Sinh viên thực hiện:} & 1. Lê Anh Khoa \\
                                    & 2. Ngô Quang Thế
    \end{tabular}
  \end{center}

  \vfill
  \begin{center}
    {\normalsize Đà Nẵng, 2026}
  \end{center}
  \vspace{0.25cm}
\end{frame}

% ==============================================================================
% SLIDE 2: MỤC LỤC NỘI DUNG (AGENDA TÁCH RIÊNG THEO YÊU CẦU)
% ==============================================================================
\begin{frame}{Mục lục Nội dung Báo cáo}
  \vspace{-0.1cm}
  \begin{columns}[t,onlytextwidth]
    \begin{column}{0.49\textwidth}
      \begin{block}{\textbf{PHẦN I: NỀN TẢNG OOP}}
        \footnotesize
        \textbf{Thuyết trình: Lê Anh Khoa}
        \begin{itemize}
          \item \textbf{Đặt vấn đề \& Phân công công việc (WBS)}
          \item \textbf{4 Trụ cột OOP trong Dart} \& Lớp bất biến (\textit{Immutable Object})
          \item \textbf{Constructor:} Named \& Redirecting Factory
          \item \textbf{Factory Constructor} \& Bảng đối sánh
          \item \textbf{Mixin:} Tái sử dụng tính năng, cơ chế \texttt{on} và lý do dùng \texttt{with}
          \item \textbf{Generics:} Reified Generics \& Bounded Type (\texttt{<T extends X>})
          \item \textbf{Sealed Class \& Pattern Matching (Dart 3)}
        \end{itemize}
      \end{block}
    \end{column}

    \begin{column}{0.49\textwidth}
      \begin{block}{\textbf{PHẦN II: NGOẠI LỆ \& KIẾN TRÚC}}
        \footnotesize
        \textbf{Thuyết trình: Ngô Quang Thế}
        \begin{itemize}
          \item \textbf{Exception vs Error:} \texttt{try-catch}, \texttt{rethrow}, \texttt{StackTrace}
          \item \textbf{Cây phân cấp Custom Exception} đa tầng
          \item \textbf{Hàm bao bọc (errorWrapper)} \& Result Pattern
          \item \textbf{Lưới an toàn toàn cục} (Global Safety Net)
          \item \textbf{Kiến trúc 3 tầng, State Management \& Chống Race Condition}
          \item \textbf{Sơ đồ Luồng dữ liệu} (Data Flow Swimlanes)
          \item \textbf{Kịch bản Thực nghiệm, Live Demo \& Tổng kết}
        \end{itemize}
      \end{block}
    \end{column}
  \end{columns}
\end{frame}

% ==============================================================================
% SLIDE 3: ĐẶT VẤN ĐỀ & BẢNG PHÂN CÔNG WBS
% ==============================================================================
\begin{frame}{Đặt vấn đề, Mục tiêu \& Bảng phân công công việc (WBS)}
  \vspace{-0.25cm}
  \textbf{1. Đặt vấn đề \& Mục tiêu đề tài:}
  \begin{itemize}
    \item \textbf{Thực trạng:} Ứng dụng di động dễ gặp sự cố crash do ngoại lệ bất đồng bộ, JSON sai định dạng, hoặc Race Condition khi người dùng thao tác dồn dập.
    \item \textbf{Mục tiêu:} Áp dụng OOP nâng cao và kỹ thuật xử lý lỗi tập trung trong Dart để xây dựng module đặt vé tự phục hồi, cách ly lỗi hoàn toàn.
  \end{itemize}

  \vspace{0.08cm}
  \textbf{2. Bảng phân công công việc chi tiết (Work Breakdown Structure):}
  \vspace{0.05cm}

  \centering
  \scriptsize
  \setlength{\tabcolsep}{4pt}
  \arrayrulecolor{white}
  \renewcommand{\arraystretch}{1.12}
  \begin{tabular}{|c|p{3.5cm}|p{5.5cm}|c|}
    \hline
    \rowcolor{TableHead}
    \makebox[0.6cm][c]{\color{white}\textbf{TT}} & \centering\color{white}\textbf{HẠNG MỤC PHỤ TRÁCH} & \centering\color{white}\textbf{FILE MÃ NGUỒN CHÍNH} & \makebox[3.0cm][c]{\color{white}\textbf{NGƯỜI THỰC HIỆN}} \rule{0pt}{2.6ex} \tabularnewline
    \hline
    \rowcolor{TableRowOdd}
    1 & \textbf{UI Design \& Safety Net} & \texttt{demo\_page.dart}, \texttt{main.dart} & Ngô Quang Thế \tabularnewline
    \hline
    \rowcolor{TableRowEven}
    2 & \textbf{Business Logic} & \texttt{booking\_service.dart}, \texttt{lab\_location.dart} & Lê Anh Khoa \tabularnewline
    \hline
    \rowcolor{TableRowOdd}
    3 & \textbf{State Management \& Utils} & \texttt{booking\_controller.dart}, \texttt{result.dart} & Ngô Quang Thế \tabularnewline
    \hline
    \rowcolor{TableRowEven}
    4 & \textbf{Database / Mock API} & \texttt{booking\_model.dart}, \texttt{app\_exception.dart} & Lê Anh Khoa \tabularnewline
    \hline
    \rowcolor{TableRowOdd}
    5 & \textbf{Reusable Logic (Mixin)} & \texttt{validation\_mixin.dart} & Lê Anh Khoa \tabularnewline
    \hline
    \rowcolor{TableRowEven}
    6 & \textbf{Báo cáo \& Slide Beamer} & Slide LaTeX, Kịch bản kiểm thử Live Demo & Cả hai thành viên \tabularnewline
    \hline
  \end{tabular}
\end{frame}

% ==============================================================================
% SLIDE 4: OOP NỀN TẢNG TRONG DART
% ==============================================================================
\begin{frame}{Lập trình Hướng đối tượng (OOP) nền tảng trong Dart}
  \vspace{-0.15cm}
  \begin{columns}[t,onlytextwidth]
    \begin{column}{0.49\textwidth}
      \footnotesize
      \textbf{1. Bốn trụ cột OOP trong Dart:}
      \begin{itemize}
        \item \textbf{Đóng gói (Encapsulation):} Ẩn giấu dữ liệu bằng tiền tố \texttt{\_} (phạm vi thư viện \textit{library-level}, không phải cấp class như Java/C++).
        \item \textbf{Kế thừa (Inheritance):} Đơn kế thừa qua \texttt{extends}, mở rộng khả năng qua \texttt{implements} (mọi class đều ngầm định là một interface).
        \item \textbf{Đa hình (Polymorphism):} Ghi đè phương thức (\texttt{@override}) và định kiểu linh hoạt qua lớp cơ sở.
        \item \textbf{Trừu tượng (Abstraction):} Khai báo qua \texttt{abstract class}, định nghĩa khung hợp đồng nghiệp vụ.
      \end{itemize}
    \end{column}

    \begin{column}{0.49\textwidth}
      \footnotesize
      \textbf{2. Thiết kế Lớp bất biến (Immutable Object):}
      \begin{itemize}
        \item \textbf{Khái niệm:} Đối tượng không thể bị thay đổi trạng thái sau khi đã được khởi tạo thành công.
        \item \textbf{Thực thi trong Dart:} Khai báo tất cả thuộc tính là \texttt{final} và sử dụng \texttt{const constructor}.
        \item \textbf{Lợi ích:} An toàn khi truyền giữa Service và Controller, dễ suy luận trong luồng bất đồng bộ (\textit{async/await}), tối ưu RAM qua chia sẻ thể hiện (\textit{Canonicalization}).
      \end{itemize}
      \vspace{0.08cm}
      \begin{block}{\scriptsize Ứng dụng trong dự án}
        \scriptsize \texttt{BookingModel} được thiết kế bất biến hoàn toàn, giúp luồng dữ liệu giữa Service và Controller luôn toàn vẹn, tránh biến đổi ngoài ý muốn.
      \end{block}
    \end{column}
  \end{columns}
\end{frame}

% ==============================================================================
% SLIDE 5: CONSTRUCTOR: NAMED & REDIRECTING FACTORY
% ==============================================================================
\begin{frame}[fragile]{Constructor: Named \& Redirecting Factory Constructor}
  \vspace{-0.2cm}
  \begin{columns}[t,onlytextwidth]
    \begin{column}{0.49\textwidth}
      \footnotesize
      \textbf{1. Named Constructor (Khởi tạo định danh):}
      \begin{itemize}
        \item Dart \textbf{không hỗ trợ Overloading constructor} cùng tên như Java/C++.
        \item Giải pháp: Đặt tên rõ nghĩa theo cú pháp \texttt{ClassName.tênRiêng()}.
        \item \textbf{Initializer List} (\texttt{: field = val}): Một trong các cách gán giá trị cho trường \texttt{final} trước khi thân hàm constructor thực thi.
      \end{itemize}

      \begin{lstlisting}[basicstyle=\ttfamily\scriptsize\color{black}, columns=flexible]
// Named constructor: initialize empty default instance
const BookingModel.empty()
    : id = 'N/A',
      amount = 0,
      type = TicketType.standard,
      labEvent = (*@'Chưa có thông tin'@*);
      \end{lstlisting}
    \end{column}

    \begin{column}{0.49\textwidth}
      \footnotesize
      \textbf{2. Redirecting Factory Constructor:}
      \begin{itemize}
        \item Dùng từ khóa \texttt{factory} kết hợp \texttt{=} trỏ trực tiếp đến constructor của lớp con tương ứng.
        \item Che giấu chi tiết cài đặt của lớp con phía sau giao diện lớp cơ sở.
      \end{itemize}

      \begin{lstlisting}[basicstyle=\ttfamily\scriptsize\color{black}, columns=flexible]
// Redirecting factory constructor: forward to subclass
sealed class Result<T> {
  const Result();
  factory Result.success(T data) 
      = Success<T>;
  factory Result.failure(AppException error) 
      = Failure<T>;
}
      \end{lstlisting}
    \end{column}
  \end{columns}
\end{frame}

% ==============================================================================
% SLIDE 6: FACTORY CONSTRUCTOR & BẢNG SO SÁNH
% ==============================================================================
\begin{frame}[fragile]{Factory Constructor \& Bảng so sánh các Constructor}
  \vspace{-0.25cm}
  \begin{columns}[t,onlytextwidth]
    \begin{column}{0.49\textwidth}
      \footnotesize
      \textbf{1. Bản chất của Factory Constructor:}
      \begin{itemize}
        \item \textbf{Không tự động tạo instance:} Bắt buộc phải có lệnh \texttt{return} trả về một đối tượng.
        \item Có thể trả về đối tượng tái sử dụng (Cache / Singleton) hoặc trả về thể hiện của một lớp con.
        \item Chứa logic nghiệp vụ, thẩm định dữ liệu và chủ động ném ngoại lệ (\texttt{throw}).
      \end{itemize}
      \begin{lstlisting}[basicstyle=\ttfamily\tiny\color{black}, columns=flexible]
// (*@BookingModel.fromJson (phiên bản minh họa rút gọn)@*)
factory BookingModel.fromJson(Map<String, dynamic> json) {
  if (json case {'id': final id, 'amount': int amount}) {
    return BookingModel(id: '$id', amount: amount, ...);
  }
  throw const DataFormatException((*@'Dữ liệu máy chủ bị lỗi'@*));
}
      \end{lstlisting}
    \end{column}

    \begin{column}{0.49\textwidth}
      \textbf{2. Bảng so sánh 3 loại Constructor:}
      \vspace{0.1cm}

      \centering
      \scriptsize
      \setlength{\tabcolsep}{2pt}
      \arrayrulecolor{white}
      \renewcommand{\arraystretch}{1.15}
      \begin{tabular}{|p{1.5cm}|c|p{1.95cm}|c|}
        \hline
        \rowcolor{TableHead}
        \centering\color{white}\textbf{TIÊU CHÍ} & \makebox[1.0cm][c]{\color{white}\textbf{NAMED}} & \makebox[2.1cm][c]{\color{white}\textbf{FACTORY}} & \makebox[1.2cm][c]{\color{white}\textbf{REDIRECT}} \rule{0pt}{2.6ex} \tabularnewline
        \hline
        \rowcolor{TableRowOdd}
        Tự tạo instance & Có & Không tự động\newline(cần \texttt{return}) & Không \tabularnewline
        \hline
        \rowcolor{TableRowEven}
        Dùng \texttt{this} / Init list & Có & Không & Không \tabularnewline
        \hline
        \rowcolor{TableRowOdd}
        Trả về lớp con & Không & Có & Có \tabularnewline
        \hline
        \rowcolor{TableRowEven}
        Chứa logic & Có (hạn chế) & Có & Không \tabularnewline
        \hline
      \end{tabular}
    \end{column}
  \end{columns}
\end{frame}

% ==============================================================================
% SLIDE 7: MIXIN TRONG DART
% ==============================================================================
\begin{frame}[fragile]{Mixin: Tái sử dụng tính năng không qua kế thừa}
  \vspace{-0.3cm}
  \begin{columns}[t,onlytextwidth]
    \begin{column}{0.49\textwidth}
      \footnotesize
      \textbf{1. Khái niệm, Linearization \& Đơn kế thừa:}
      \begin{itemize}
        \item \textbf{Bản chất:} Tái sử dụng mã nguồn giữa các class không cùng cây phả hệ.
        \item \textbf{Cơ chế Linearization:} Dart giải quyết xung đột phương thức bằng \textbf{Linearization} (thứ tự khai báo từ trái sang phải: mixin sau ghi đè mixin trước), giải quyết xung đột hình thoi (\textit{Diamond Problem}).
        \item \textbf{Từ khóa \texttt{on}:} Giới hạn kiểu được phép dùng mixin (\texttt{mixin M on SuperClass}), cho phép gọi \texttt{super.method()}.
        \item \textbf{Lý do dùng \texttt{with}:} Dart là ngôn ngữ \textbf{đơn kế thừa}. \texttt{BookingController} đã \texttt{extends ChangeNotifier} nên \textbf{bắt buộc dùng \texttt{with}} để nhúng \texttt{ValidationAndLogMixin}.
      \end{itemize}
    \end{column}

    \begin{column}{0.49\textwidth}
      \footnotesize
      \textbf{2. Bảng đối sánh \& Code minh họa trong dự án:}
      \vspace{0.04cm}

      \centering
      \scriptsize
      \setlength{\tabcolsep}{2pt}
      \arrayrulecolor{white}
      \renewcommand{\arraystretch}{1.1}
      \begin{tabular}{|>{\raggedright\arraybackslash}p{2.1cm}|c|c|c|}
        \hline
        \rowcolor{TableHead}
        \centering\color{white}\textbf{ĐẶC TRƯNG} & \makebox[1.2cm][c]{\color{white}\textbf{EXTENDS}} & \makebox[1.2cm][c]{\color{white}\textbf{IMPL}} & \makebox[1.2cm][c]{\color{white}\textbf{WITH}} \rule{0pt}{2.2ex} \tabularnewline
        \hline
        \rowcolor{TableRowOdd}
        Bản chất & "Là một" & "Hợp đồng" & "Thêm hành vi" \tabularnewline
        \hline
        \rowcolor{TableRowEven}
        Số lượng & Duy nhất 1 & Nhiều & Nhiều \tabularnewline
        \hline
        \rowcolor{TableRowOdd}
        Thừa hưởng & Có & Không & Có \tabularnewline
        \hline
        \rowcolor{TableRowEven}
        Có constructor riêng & Có & Có & Không \tabularnewline
        \hline
      \end{tabular}

      \vspace{0.04cm}
      \begin{lstlisting}[basicstyle=\ttfamily\scriptsize\color{black}, columns=flexible]
mixin ValidationAndLogMixin {
  void validateTicketCount(int count) {
    if (count <= 0) {
      throw const BusinessLogicException(
          (*@'Số lượng vé phải lớn hơn 0'@*));
    }
  }
  void logAction(String a) => print('[LOG] $a');
}
      \end{lstlisting}
    \end{column}
  \end{columns}
\end{frame}

% ==============================================================================
% SLIDE 8: GENERICS CƠ BẢN TRONG DART
% ==============================================================================
\begin{frame}[fragile]{Generics: Đảm bảo an toàn kiểu \& Reified Generics}
  \vspace{-0.3cm}
  \begin{columns}[t,onlytextwidth]
    \begin{column}{0.49\textwidth}
      \footnotesize
      \textbf{1. Mục đích và Ý nghĩa:}
      \begin{itemize}
        \item Cho phép viết code tổng quát hoạt động trên nhiều kiểu dữ liệu mà vẫn đảm bảo \textbf{an toàn kiểu lúc biên dịch} (\textit{Compile-time Type Safety}).
        \item Loại bỏ thao tác ép kiểu thủ công không an toàn (\mbox{\texttt{as Type}}), phòng ngừa phát sinh lỗi \texttt{TypeError} khi chạy.
      \end{itemize}

      \vspace{0.08cm}
      \textbf{2. Cơ chế Reified Generics trong Dart:}
      \begin{itemize}
        \item Khác Java (\textit{Type Erasure}), Dart \textbf{lưu thông tin kiểu generic cả trong lúc chạy} (\textit{Runtime}).
        \item Cho phép kiểm tra kiểu chính xác lúc thực thi: \texttt{result is Result<BookingModel>}.
      \end{itemize}
    \end{column}

    \begin{column}{0.49\textwidth}
      \footnotesize
      \textbf{3. Ràng buộc Generics (\mbox{\texttt{<T extends X>}}):}
      \begin{itemize}
        \item \textbf{Bounded Type:} Giới hạn kiểu tham số nhận lớp con của kiểu định sẵn.
        \item \textbf{Non-nullable Bound:} Với Null Safety, \texttt{<T>} mặc định là \texttt{Object?}. Viết \mbox{\texttt{<T extends Object>}} nhằm \textbf{loại trừ triệt để kiểu null}.
      \end{itemize}
      \begin{lstlisting}[basicstyle=\ttfamily\scriptsize\color{black}, columns=flexible]
// (*@Minh họa Bounded Type: chỉ nhận lớp con của AppException@*)
class ErrorHandler<E extends AppException> {
  final E error;
  const ErrorHandler(this.error);
  void log() => print('[${error.code}] ${error.message}');
}
      \end{lstlisting}
      \vspace{0.06cm}
      {\scriptsize\textit{$\rightarrow$ Ứng dụng trong dự án: \texttt{Result<T>}, \texttt{errorWrapper<T>}}}
    \end{column}
  \end{columns}
\end{frame}

% ==============================================================================
% SLIDE 9: SEALED CLASS & PATTERN MATCHING (DART 3)
% ==============================================================================
\begin{frame}[fragile]{Sealed Class \& Pattern Matching (Dart 3)}
  \begin{columns}[t,onlytextwidth]
    \begin{column}{0.49\textwidth}
      \textbf{1. Tính năng Sealed Class trong Dart 3:}
      \begin{itemize}
        \item \texttt{sealed class Result<T>}: Khóa cây kế thừa, chỉ cho phép các lớp con (\texttt{Success<T>}, \texttt{Failure<T>}) khai báo trong cùng một thư viện (\textit{same library}).
        \item \textbf{Exhaustive Pattern Matching:} Trình biên dịch bắt buộc biểu thức \texttt{switch} phải kiểm tra đủ mọi trường hợp; báo lỗi ngay nếu bỏ sót nhánh.
        \item \textbf{Loại bỏ trạng thái nhập nhằng:} Đảm bảo không bao giờ tồn tại trạng thái không hợp lệ (\textit{"vừa có dữ liệu vừa có lỗi"}).
      \end{itemize}
    \end{column}

    \begin{column}{0.49\textwidth}
      \textbf{2. Cú pháp Destructuring tinh gọn:}
      \begin{lstlisting}[basicstyle=\ttfamily\scriptsize\color{black}, columns=flexible]
// Exhaustive pattern matching on sealed Result<T>
switch (result) {
  case Success(:final data):
    isLastActionSuccess = true;
    lastBookedTicket = data;
    location.deductSeats(type, data.amount);
    
  case Failure(:final error):
    isLastActionSuccess = false;
    statusMessage = error.message;
}
      \end{lstlisting}
    \end{column}
  \end{columns}
  \note{Lê Anh Khoa hoàn thành trình bày Phần I về OOP nền tảng. Chuyển giao phần trình bày cho Ngô Quang Thế tiếp tục với Phần II về Xử lý ngoại lệ và Kiến trúc.}
\end{frame}

% ==============================================================================
% SLIDE 10: XỬ LÝ NGOẠI LỆ TRONG DART
% ==============================================================================
\begin{frame}[fragile]{Xử lý Ngoại lệ trong Dart: Exception vs Error}
  \vspace{-0.35cm}
  \begin{columns}[t,onlytextwidth]
    \begin{column}{0.49\textwidth}
      \footnotesize
      \textbf{1. Phân biệt Exception và Error:}
      \begin{itemize}
        \item \textbf{Exception:} Ngoại lệ dự đoán trước trong runtime (mạng, JSON). Ứng dụng \textbf{phải bắt và tự phục hồi}.
        \item \textbf{Error:} Lỗi lập trình (\texttt{TypeError}, \texttt{RangeError}). Về nguyên tắc không nên bắt mà cần sửa code \textit{(chỉ bắt tại biên để chống crash \& log)}.
      \end{itemize}

      \vspace{0.05cm}
      \textbf{2. Từ khóa \texttt{rethrow} \& \texttt{StackTrace}:}
      \begin{itemize}
        \item \texttt{catch (e, stackTrace)}: Bắt đối tượng ngoại lệ kèm vết ngăn xếp chi tiết.
        \item \texttt{rethrow}: Đẩy lỗi lên trên mà vẫn \textbf{bảo toàn Stack Trace} gốc.
      \end{itemize}

      \vspace{0.05cm}
      \textbf{3. Hai quy tắc vàng khi bắt ngoại lệ:}
      \begin{itemize}
        \item Khối \texttt{on} cụ thể (\texttt{on AppException}) \textbf{phải đặt trước} \texttt{catch} chung để tránh bị nuốt mất ngoại lệ.
        \item Khối \texttt{finally} \textbf{luôn thực thi} dọn tài nguyên kể cả khi khối lệnh trước đã gọi \texttt{return}.
      \end{itemize}
    \end{column}

    \begin{column}{0.49\textwidth}
      \footnotesize\textbf{4. Cấu trúc đầy đủ try-on-catch-finally:}
      \begin{lstlisting}[basicstyle=\ttfamily\scriptsize\color{black}, columns=flexible]
try {
  final res = await apiCall();
} on AppException catch (e) {
  // Catch typed domain exception
  logError(e.message);
} catch (e, stackTrace) {
  // Catch unhandled & inspect trace
  logError('Unhandled: $e\n$stackTrace');
} finally {
  // (*@Luôn chạy kể cả khi try đã return!@*)
  closeConnection();
}
      \end{lstlisting}
    \end{column}
  \end{columns}
\end{frame}

% ==============================================================================
% SLIDE 11: CÂY PHÂN CẤP CUSTOM EXCEPTION
% ==============================================================================
\begin{frame}[fragile]{Cây phân cấp Custom Exception (\texttt{AppException})}
  \vspace{-0.35cm}
  \begin{columns}[t,onlytextwidth]
    \begin{column}{0.48\textwidth}
      \textbf{1. Triết lý phân cấp ngoại lệ tùy biến:}
      \begin{itemize}
        \item Không ném \texttt{Exception} thô, luôn dùng lớp có \texttt{code} và \texttt{message} chuẩn hóa.
        \item Xây dựng lớp trừu tượng \texttt{AppException} triển khai (\texttt{implements}) interface \texttt{Exception}.
        \item Mỗi phân lớp phản ánh chính xác một nhóm lỗi nghiệp vụ trong ứng dụng.
      \end{itemize}

      \vspace{0.05cm}
      \begin{block}{\scriptsize 4 Phân nhóm ngoại lệ chuẩn hóa}
        \tiny
        $\bullet$ \textbf{NET\_ERR:} Mất kết nối mạng / máy chủ.\\
        $\bullet$ \textbf{PARSE\_ERR:} Lỗi giải mã dữ liệu JSON.\\
        $\bullet$ \textbf{BIZ\_ERR:} Vi phạm quy định số lượng vé.\\
        $\bullet$ \textbf{UNKNOWN\_ERR:} Sự cố hệ thống bất ngờ.
      \end{block}
    \end{column}

    \begin{column}{0.50\textwidth}
      \textbf{2. Triển khai kế thừa với Super Parameters (Dart 2.17+):}
      \begin{lstlisting}[basicstyle=\ttfamily\tiny\color{black}, columns=flexible]
abstract class AppException implements Exception {
  final String message, code;
  const AppException(this.message, {required this.code});
  @override
  String toString() => '[$code] $message';
}

class NetworkException extends AppException {
  const NetworkException(super.message)
      : super(code: 'NET_ERR');
}
class DataFormatException extends AppException {
  const DataFormatException(super.message)
      : super(code: 'PARSE_ERR');
}
class BusinessLogicException extends AppException {
  const BusinessLogicException(super.message)
      : super(code: 'BIZ_ERR');
}
class UnknownException extends AppException {
  const UnknownException(super.message)
      : super(code: 'UNKNOWN_ERR');
}
      \end{lstlisting}
    \end{column}
  \end{columns}
\end{frame}

% ==============================================================================
% SLIDE 12: HÀM BAO BỌC (ERROR WRAPPER) & RESULT PATTERN
% ==============================================================================
\begin{frame}[fragile]{Hàm bao bọc (Error Wrapper) \& Result Pattern}
  \vspace{-0.35cm}
  \begin{columns}[t,onlytextwidth]
    \begin{column}{0.49\textwidth}
      \footnotesize
      \textbf{1. Cơ chế đóng gói an toàn với \texttt{errorWrapper}:}
      \begin{itemize}
        \item \textbf{Cổng Fail-safe ở tầng biên:} \texttt{catch (e, st)} bắt trọn mọi lỗi (kể cả Error) để tránh sập app, đồng thời ghi log chi tiết kèm \texttt{StackTrace}.
        \item \textbf{Chủ ý của \texttt{finally}:} Luôn thực thi dọn tài nguyên và log kết thúc phiên dù thành công hay ném lỗi.
      \end{itemize}

      \begin{lstlisting}[basicstyle=\ttfamily\tiny\color{black}, columns=flexible]
// (*@errorWrapper (phiên bản minh họa rút gọn)@*)
Future<Result<T>> errorWrapper<T>(
    Future<T> Function() action) async {
  try {
    return Result.success(await action());
  } on AppException catch (e) {
    _recordError(e.toString()); // Log technical
    return Result.failure(e);   // Domain typed
  } catch (e, stackTrace) {
    _recordError('Unhandled: $e\n$stackTrace');
    return Result.failure(const UnknownException(
        (*@'Đã có lỗi xảy ra trong quá trình xử lý, vui lòng thử lại'@*)));
  } finally {
    // (*@Chủ ý: luôn chạy đóng phiên kể cả khi có lỗi@*)
    (*@logInfo('Wrapper', 'Hoàn tất phiên xử lý');@*)
  }
}
      \end{lstlisting}
    \end{column}

    \begin{column}{0.49\textwidth}
      \footnotesize
      \textbf{2. Ma trận điều phối \& thông báo thân thiện:}
      \vspace{0.08cm}

      \centering
      \scriptsize
      \setlength{\tabcolsep}{2.5pt}
      \arrayrulecolor{white}
      \renewcommand{\arraystretch}{1.15}
      \begin{tabular}{|p{1.8cm}|c|p{2.95cm}|}
        \hline
        \rowcolor{TableHead}
        \centering\color{white}\textbf{NGUYÊN NHÂN} & \makebox[1.3cm][c]{\color{white}\textbf{MÃ LỖI}} & \centering\color{white}\textbf{THÔNG BÁO UI} \rule{0pt}{2.4ex} \tabularnewline
        \hline
        \rowcolor{TableRowOdd}
        Số vé $\le 0$ / Hết chỗ & \texttt{BIZ\_ERR} & \textit{"Số lượng vé phải lớn hơn 0"}\newline\textit{"Không đủ vé tại phòng lab..."} \tabularnewline
        \hline
        \rowcolor{TableRowEven}
        Mất mạng / Server sập & \texttt{NET\_ERR} & Mất kết nối tới máy chủ \tabularnewline
        \hline
        \rowcolor{TableRowOdd}
        Payload JSON hỏng & \texttt{PARSE\_ERR} & Dữ liệu máy chủ bị lỗi \tabularnewline
        \hline
        \rowcolor{TableRowEven}
        Lỗi bất ngờ khác & \texttt{UNKNOWN\_ERR} & Đã có lỗi xảy ra trong quá trình xử lý, vui lòng thử lại \tabularnewline
        \hline
      \end{tabular}
    \end{column}
  \end{columns}
\end{frame}

% ==============================================================================
% SLIDE 13: LƯỚI AN TOÀN TOÀN CỤC (GLOBAL SAFETY NET)
% ==============================================================================
\begin{frame}[fragile]{Lưới an toàn toàn cục (Global Safety Net)}
  \vspace{-0.3cm}
  \begin{columns}[t,onlytextwidth]
    \begin{column}{0.49\textwidth}
      \footnotesize
      \textbf{1. Hệ thống phòng thủ 2 lớp (Defense-in-Depth):}
      \begin{itemize}
        \item \textbf{Lớp 1 (Framework):} \texttt{FlutterError.onError} bắt lỗi tầng framework (build/render widget, layout, gestures...).
        \item \textbf{Lớp 2 (Platform):} \texttt{PlatformDispatcher.onError} bắt ngoại lệ bất đồng bộ lọt ra ngoài luồng.
        \item Trả về \texttt{true} để báo đã xử lý, ngăn lan truyền mặc định; lỗi được gom vào Console và hiển thị an toàn trên UI.
      \end{itemize}

      \vspace{0.02cm}
      \begin{block}{\scriptsize Đồng bộ lỗi tức thì}
        \scriptsize Lỗi đẩy vào \texttt{errorLogs} của Controller và cập nhật lên Console Terminal màu đen trên giao diện.
      \end{block}
    \end{column}

    \begin{column}{0.49\textwidth}
      \textbf{2. Cài đặt thực tế trong \texttt{main.dart}:}
      \begin{lstlisting}[basicstyle=\ttfamily\scriptsize\color{black}, columns=flexible]
void main() {
  WidgetsFlutterBinding.ensureInitialized();
  final controller = BookingController();

  // Layer 1: Catch framework UI render errors
  FlutterError.onError = (details) {
    FlutterError.presentError(details);
    controller.recordExternalError(
      'FlutterFramework', details.exception);
  };

  // Layer 2: Catch unhandled async exceptions
  PlatformDispatcher.instance.onError = 
      (error, stack) {
    controller.recordExternalError(
      'PlatformDispatcher', error);
    return true; // Prevent app crash
  };

  runApp(MyApp(controller: controller));
}
      \end{lstlisting}
    \end{column}
  \end{columns}
\end{frame}

% ==============================================================================
% SLIDE 14: KIẾN TRÚC GIẢI PHÁP & CHỐNG RACE CONDITION
% ==============================================================================
\begin{frame}{Kiến trúc Phân tầng, State Management \& Chống Race Condition}
  \vspace{-0.35cm}
  \begin{columns}[c,onlytextwidth]
    \begin{column}{0.50\textwidth}
      \scriptsize
      \textbf{1. Phân tầng kiến trúc rạch ròi:}
      \begin{itemize}
        \setlength{\itemsep}{1pt}
        \item \textbf{Presentation Layer:} \texttt{DemoPage} $\leftrightarrow$ \texttt{ListenableBuilder} $\leftrightarrow$ \texttt{Controller}.
        \item \textbf{Reusable Mixin:} \texttt{ValidationAndLogMixin} (dùng chung Controller \& Service).
        \item \textbf{Data Layer:} \texttt{BookingService} (\texttt{errorWrapper}), \texttt{Model}, \texttt{Location}.
        \item \textbf{Core Layer:} \texttt{Result<T>} (Sealed Class), \texttt{AppException} (Cây ngoại lệ).
      \end{itemize}

      \vspace{0.04cm}
      \textbf{2. Bất đồng bộ \& Quản lý trạng thái:}
      \begin{itemize}
        \setlength{\itemsep}{1pt}
        \item Áp dụng mô hình bất đồng bộ (\texttt{Future}, \texttt{async/await}, Event Loop).
        \item \texttt{BookingController} phát thông báo qua \texttt{notifyListeners()}; giao diện dùng \texttt{ListenableBuilder} tái tạo chỉ các widget cần thiết.
      \end{itemize}

      \vspace{0.04cm}
      \textbf{3. Chống Race Condition (Bấm liên tiếp):}
      \begin{itemize}
        \setlength{\itemsep}{1pt}
        \item Chắn đầu tác vụ: \texttt{if (isLoading) return;}.
        \item Khóa tương tác UI: \texttt{onPressed: isLoading ? null : \_submitBooking}.
      \end{itemize}
    \end{column}

    \begin{column}{0.48\textwidth}
      \centering
      \begin{tikzpicture}[
        node distance=0.28cm,
        box/.style={draw=black, fill=white, thick, rounded corners=3pt, minimum width=6.1cm, minimum height=0.48cm, font=\tiny\bfseries, align=center},
        safety/.style={draw=black, fill=white, thick, dashed, rounded corners=3pt, minimum width=6.1cm, minimum height=0.45cm, font=\tiny\bfseries, align=center},
        mixin/.style={draw=black, fill=white, thick, rounded corners=3pt, minimum width=6.1cm, minimum height=0.45cm, font=\tiny\bfseries, align=center},
        arrow/.style={thick, ->, >=Stealth, draw=black}
      ]
        \node[safety] (main) {GLOBAL SAFETY NET (main.dart)\\FlutterError + PlatformDispatcher};
        \node[box, below=of main] (pres) {PRESENTATION LAYER\\DemoPage $\leftrightarrow$ ListenableBuilder $\leftrightarrow$ Controller};
        \node[mixin, below=of pres] (mixin) {REUSABLE MIXIN\\ValidationAndLogMixin (Log, Time, Validate)};
        \node[box, below=of mixin] (data) {DATA LAYER\\BookingService (errorWrapper) $\leftrightarrow$ Model / Location};
        \node[box, below=of data] (core) {CORE LAYER\\Result<T> (Sealed Class) $\leftrightarrow$ AppException Hierarchy};

        \draw[arrow] (main) -- node[right, font=\tiny] {bảo vệ UI} (pres);
        \draw[arrow] (pres) -- node[right, font=\tiny] {with mixin} (mixin);
        \draw[arrow] (mixin) -- node[right, font=\tiny] {with mixin} (data);
        \draw[arrow] (data) -- node[right, font=\tiny] {dùng kiểu} (core);
      \end{tikzpicture}
    \end{column}
  \end{columns}
\end{frame}

% ==============================================================================
% SLIDE 15: SƠ ĐỒ LUỒNG DỮ LIỆU (DATA FLOW PIPELINE DẠNG LÀN DỌC)
% ==============================================================================
\begin{frame}{Sơ đồ Luồng dữ liệu (Data Flow Pipeline dạng Làn dọc)}
  \vspace{-0.3cm}
  \centering
  \resizebox{0.99\textwidth}{!}{%
  \begin{tikzpicture}[
    lane/.style={draw=black!40, dashed, fill=white, rounded corners=4pt},
    stepbox/.style={draw=black, fill=white, thick, rounded corners=3pt, minimum width=2.65cm, minimum height=0.68cm, inner sep=3.5pt, font=\scriptsize\bfseries, align=center},
    errbox/.style={draw=black, dashed, fill=white, thick, rounded corners=3pt, minimum width=2.85cm, minimum height=0.68cm, inner sep=3.5pt, font=\scriptsize\bfseries, align=center},
    succbox/.style={draw=black, double, fill=white, thick, rounded corners=3pt, minimum width=2.85cm, minimum height=0.68cm, inner sep=3.5pt, font=\scriptsize\bfseries, align=center},
    arrow/.style={thick, ->, >=Stealth, draw=black},
    retarrow/.style={thick, dashed, ->, >=Stealth, draw=black}
  ]
    % 4 Làn dọc (Swimlanes) - Thiết kế mở rộng, thoáng đãng, chuẩn tọa độ
    \draw[lane] (0.0, 0) rectangle (3.0, -6.2);
    \draw[lane] (3.3, 0) rectangle (6.5, -6.2);
    \draw[lane] (6.8, 0) rectangle (11.5, -6.2);
    \draw[lane] (11.8, 0) rectangle (15.6, -6.2);

    % Tiêu đề làn
    \node[font=\small\bfseries\color{black}] at (1.50, 0.25) {1. UI (DemoPage)};
    \node[font=\small\bfseries\color{black}] at (4.90, 0.25) {2. Controller};
    \node[font=\small\bfseries\color{black}] at (9.15, 0.25) {3. Service + Wrapper};
    \node[font=\small\bfseries\color{black}] at (13.70, 0.25) {4. Model \& Location};

    % Hàng 1 (y = -0.75): Khởi tạo tác vụ
    \node[stepbox] (s1) at (1.50, -0.75) {1. Bấm "Đặt vé"\\\relax (isLoading = true)};
    \node[stepbox] (s2) at (4.90, -0.75) {2. bookTicket()\\\relax Khóa tác vụ trùng};
    \node[stepbox, minimum width=2.9cm] (s3) at (9.15, -0.75) {3. errorWrapper<T>\\\relax Mở khối try-catch};

    % Hàng 2 (y = -2.00): Kiểm tra hợp lệ dữ liệu & Nhánh lỗi nghiệp vụ
    \node[stepbox, minimum width=3.4cm] (s4) at (13.70, -2.00) {4. validateTicketCount() (Mixin)\\\relax checkAvailability() (Location)};
    \node[errbox] (e1) at (9.15, -2.00) {Lỗi vé / Hết chỗ\\\relax [BIZ\_ERR]};

    % Hàng 3 (y = -3.25): Giả lập gọi API, Parse JSON & Nhánh lỗi hạ tầng
    \node[stepbox, minimum width=3.2cm] (s5) at (13.70, -3.25) {5. Giả lập API / JSON\\\relax BookingModel.fromJson()};
    \node[errbox, minimum width=3.1cm] (e2) at (9.15, -3.25) {Mất mạng / JSON hỏng\\\relax [NET/PARSE\_ERR]};

    % Hàng 4 (y = -4.45): Kết quả Thất bại (Failure)
    \node[errbox, minimum width=3.3cm] (fail) at (9.15, -4.45) {Result.failure(error)\\\relax Ghi log vào Console};
    \node[errbox] (c_fail) at (4.90, -4.45) {6. Failure: Gán thông báo\\\relax statusMessage = err.msg};

    % Hàng 5 (y = -5.55): Kết quả Thành công (Success)
    \node[succbox, minimum width=3.3cm] (succ) at (9.15, -5.55) {Result.success(ticket)\\\relax Đóng gói vé thành công};
    \node[succbox] (c_succ) at (4.90, -5.55) {6. Success: Trừ ghế\\\relax Lưu vé vào lịch sử};

    % Cột 1: Hoàn tất tác vụ, cập nhật State & Giao diện (y = -5.00)
    \node[stepbox, minimum height=0.95cm] (s_ui) at (1.50, -5.00) {7. notifyListeners()\\\relax Render UI an toàn\\\relax (isLoading = false)};

    % Luồng khởi đầu: UI -> Controller -> Service
    \draw[arrow] (s1) -- (s2);
    \draw[arrow] (s2) -- (s3);

    % Luồng thực thi: Service -> Model (Validation) -> Simulation/Parse
    \draw[arrow] (s3.east) -| node[above, pos=0.35, font=\scriptsize\bfseries\color{black}] {Thực thi} (s4.north);
    \draw[arrow] (s4.south) -- node[right, font=\scriptsize\bfseries\color{black}] {Hợp lệ} (s5.north);

    % Luồng ném ngoại lệ: Model -> Service (2 nhánh độc lập hoàn toàn, rộng rãi)
    \draw[arrow] (s4.west) -- node[above, font=\scriptsize\bfseries\color{black}] {Vi phạm} (e1.east);
    \draw[arrow] (s5.west) -- node[above, font=\scriptsize\bfseries\color{black}] {Sự cố} (e2.east);

    % Hai nhánh lỗi đổ về Result.failure (tách rời cân đối hai bên, không cắt nhau)
    \draw[arrow] (e1.west) -- ++(-0.6, 0) |- ([xshift=-0.8cm]fail.north);
    \draw[arrow] (e2.south) -- (fail.north);

    % Luồng thành công đóng gói vé: Model -> Service (Result.success)
    \draw[arrow] (s5.south) -- ++(0, -1.0) node[right=2pt, font=\scriptsize\bfseries\color{black}] {Thành công} |- (succ.east);

    % Luồng trả kết quả từ Service về Controller (2 đường song song độc lập)
    \draw[retarrow] (fail.west) -- (c_fail.east);
    \draw[retarrow] (succ.west) -- (c_succ.east);

    % Cập nhật giao diện UI (quy tụ về s_ui)
    \draw[retarrow] (c_fail.west) -- ([yshift=7pt]s_ui.east);
    \draw[retarrow] (c_succ.west) -- ([yshift=-7pt]s_ui.east);
  \end{tikzpicture}%
  }

  \vspace{0.08cm}
  \begin{center}
    \scriptsize \textbf{Nguyên tắc cốt lõi:} Mọi nhánh lỗi hay thành công đều được chuẩn hóa qua \texttt{Result<T>} trả về Controller; tầng UI hoàn toàn sạch sẽ, không chứa \texttt{try-catch}.
  \end{center}
\end{frame}

% ==============================================================================
% SLIDE 16: KỊCH BẢN THỰC NGHIỆM & LIVE DEMO
% ==============================================================================
\begin{frame}{Kịch bản Thực nghiệm \& Bảng ca kiểm thử Live Demo}
  \vspace{-0.25cm}
  \textbf{Ma trận kiểm thử nhanh tích hợp trên giao diện \& ánh xạ kỹ thuật OOP:}
  \vspace{0.06cm}

  \centering
  \scriptsize
  \setlength{\tabcolsep}{3.5pt}
  \arrayrulecolor{white}
  \renewcommand{\arraystretch}{1.12}
  \begin{tabular}{|p{2.6cm}|>{\raggedright\arraybackslash}p{3.2cm}|>{\raggedright\arraybackslash}p{5.1cm}|c|}
    \hline
    \rowcolor{TableHead}
    \centering\color{white}\textbf{NÚT KIỂM THỬ} & \centering\color{white}\textbf{KỸ THUẬT OOP ÁP DỤNG} & \centering\color{white}\textbf{KẾT QUẢ GIAO DIỆN (UI BANNER)} & \makebox[1.8cm][c]{\color{white}\textbf{MÃ LỖI}} \rule{0pt}{2.3ex} \tabularnewline
    \hline
    \rowcolor{TableRowOdd}
    \textbf{Đặt 2 vé chuẩn} & Factory fromJson + Result.success (Sealed) & Đặt thành công 2 vé, khấu trừ ghế trống phòng lab & --- \tabularnewline
    \hline
    \rowcolor{TableRowEven}
    \textbf{Số vé âm (-1)} & Mixin: \texttt{validateTicketCount()} & Báo lỗi: \textit{"Số lượng vé phải lớn hơn 0"} & \texttt{[BIZ\_ERR]} \tabularnewline
    \hline
    \rowcolor{TableRowOdd}
    \textbf{Vượt vé chuẩn (50 vé)} & LabLocation: \texttt{checkAvailability()} & Báo lỗi: \textit{"Không đủ vé Tiêu chuẩn tại [Hội thảo...]"} & \texttt{[BIZ\_ERR]} \tabularnewline
    \hline
    \rowcolor{TableRowEven}
    \textbf{Vượt vé VIP (16 vé)} & LabLocation: \texttt{checkAvailability()} & Báo lỗi: \textit{"Không đủ vé VIP tại [Hội thảo...]"} & \texttt{[BIZ\_ERR]} \tabularnewline
    \hline
    \rowcolor{TableRowOdd}
    \textbf{15 vé VIP (Mất mạng)} & \texttt{errorWrapper<T>} bắt ngoại lệ & Báo lỗi: \textit{"Mất kết nối tới máy chủ xác thực đặc quyền..."} & \texttt{[NET\_ERR]} \tabularnewline
    \hline
    \rowcolor{TableRowEven}
    \textbf{Lab 404 (Lỗi JSON)} & Factory Constructor: \texttt{fromJson} & Báo lỗi: \textit{"Dữ liệu máy chủ bị lỗi"} & \texttt{[PARSE\_ERR]} \tabularnewline
    \hline
  \end{tabular}

  \vspace{0.08cm}
  \begin{block}{\scriptsize Nguyên tắc liên kết kịch bản kiểm thử với thiết kế OOP}
    \scriptsize
    $\bullet$ \textbf{Kiểm soát tại nguồn:} Lỗi bị chặn ngay tại tầng phát sinh (\texttt{validateTicketCount()}, \texttt{checkAvailability()}, \texttt{fromJson}).\\
    $\bullet$ \textbf{Cách ly lỗi tập trung:} Cổng \texttt{errorWrapper<T>} chuẩn hóa ngoại lệ thành \texttt{Result.failure}, hiển thị an toàn trên UI.
  \end{block}
\end{frame}

% ==============================================================================
% SLIDE 17: MINH CHỨNG THỰC NGHIỆM - TOÀN BỘ GIAO DIỆN ỨNG DỤNG
% ==============================================================================
\begin{frame}{Minh chứng Thực nghiệm: Toàn bộ Giao diện Ứng dụng}
  \vspace{-0.45cm}
  \textbf{Giao diện thực tế ứng dụng Smart Booking System \& Nhật ký gom lỗi Console:}
  \par\vspace{0.04cm}
  \centering
  \fboxsep=0pt\fboxrule=0.6pt%
  \fbox{\includegraphics[width=0.935\textwidth]{app_full_screen.png}}
  \par\vspace{0.04cm}
  {\scriptsize\textit{Ghi chú: Banner trên cùng phản hồi lỗi thân thiện; Console dưới cùng ghi nhận mã lỗi kỹ thuật. Mỗi sự kiện có kho vé riêng: Hội thảo SoC \& FPGA còn 33/35 sau khi đặt 2 vé; Lab 404 vẫn 35/35.}}
\end{frame}

% ==============================================================================
% SLIDE 18: TỔNG KẾT ĐỀ TÀI & HƯỚNG PHÁT TRIỂN (TÁCH RIÊNG THEO YÊU CẦU)
% ==============================================================================
\begin{frame}{Tổng kết Đề tài \& Hướng phát triển Tương lai}
  \vspace{-0.15cm}
  \begin{columns}[t,onlytextwidth]
    \begin{column}{0.49\textwidth}
      \textbf{1. Bốn kết quả cốt lõi đạt được:}
      \begin{itemize}
        \item \textbf{Named \& Factory Constructor:} Khởi tạo linh hoạt, chuyển đổi an toàn từ JSON bằng Pattern Matching Dart 3.
        \item \textbf{Mixin:} Tái sử dụng tính năng validate và log thời gian giữa các tầng mà không cần kế thừa phức tạp.
        \item \textbf{Generics + Sealed Class:} Một wrapper duy nhất cho mọi kiểu dữ liệu, ép xử lý triệt để nhánh kết quả.
        \item \textbf{Custom Exception + Result + Safety Net:} Ứng dụng không crash, thông báo thân thiện, gom lỗi tập trung.
      \end{itemize}
    \end{column}

    \begin{column}{0.49\textwidth}
      \textbf{2. Hạn chế \& Hướng phát triển tương lai:}
      \begin{itemize}
        \item \textbf{Hạn chế hiện tại:} Dữ liệu đang được giả lập trên RAM, chưa lưu trữ bền vững.
        \item \textbf{Hướng phát triển:}
          \begin{itemize}
            \item Kết nối RESTful API thật (bắt \texttt{SocketException}, \texttt{TimeoutException}).
            \item Lưu trữ dữ liệu cục bộ với SQLite hoặc Hive.
            \item Mở rộng State Management chuyên nghiệp với Bloc / Riverpod.
            \item Xây dựng hệ thống kiểm thử tự động (Unit / Widget Test) toàn diện.
          \end{itemize}
      \end{itemize}
    \end{column}
  \end{columns}
\end{frame}

% ==============================================================================
% SLIDE 19: LỜI CẢM ƠN & PHIÊN HỎI ĐÁP (Q&A TÁCH RIÊNG THEO YÊU CẦU)
% ==============================================================================
\begin{frame}[plain]
  \vspace{1.2cm}
  \begin{center}
    {\Huge \textbf{CẢM ƠN THẦY CÔ}}\\[0.25cm]
    {\Huge \textbf{VÀ CÁC BẠN ĐÃ LẮNG NGHE!}}\\[0.6cm]
    
    {\Large \textbf{PHIÊN THẢO LUẬN \& HỎI ĐÁP (Q\&A)}}\\[0.3cm]
    {\normalsize Nhóm sẵn sàng lắng nghe mọi câu hỏi, phản biện và góp ý từ quý Thầy Cô!}
  \end{center}

  \vspace{0.8cm}
  \begin{center}
    \small
    \textbf{Sinh viên thực hiện:}\\[0.1cm]
    1. Lê Anh Khoa \quad -- \quad 2. Ngô Quang Thế\\[0.15cm]
    \textit{Khoa Điện tử và Trí tuệ Nhân tạo -- Trường Đại học Bách Khoa}
  \end{center}
\end{frame}

\end{document}
